\paragraph{Timing in simultaneous translation and interpreting.}
EVS and related measures provide reproducible descriptions of temporal offset when source--target linguistic anchors are explicitly aligned; automated EVS work makes the alignment choice part of the measurement problem~\citep{guo2024manual}. EVS also varies with text-specific factors, interpreter-specific factors, and individual variation~\citep{janikowski2025evs}, alongside segmentation, reformulation, and production decisions central to cognitive accounts of SI~\citep{gile1995effort,seeber2011cognitive}. In simultaneous machine translation, wait-$k$ and prefix-to-prefix policies make latency controllable during generation~\citep{ma2019stacl}, and extensions to simultaneous speech translation analyze latency--quality trade-offs in end-to-end settings~\citep{ma2020simulst}. Average Token Delay (ATD) further argues that latency metrics should account for output duration, not only the start timing of partial translations~\citep{kano2023atd}. Our timing feature instead uses segment-onset delay: the archived onset difference for a paired source/interpreted segment, without word-level anchors or manual alignment validation. It serves as a baseline and input feature for predicting professional promptness ratings.

\paragraph{Cognitive accounts of SI.}
Gile's Effort Model characterizes SI as concurrent listening/analysis, production, and memory demands~\citep{gile1995effort}. Resource-allocation accounts further model how these demands interact under cognitive load~\citep{seeber2011cognitive,seeber2012cognitive}. Such accounts motivate the possibility that timing, fluency, and content quality co-vary in professional judgments. We test whether quality-associated signals help recover professional promptness labels on held-out speeches.

\paragraph{Automatic quality estimation.}
COMET introduced learned neural evaluation from source and translation representations~\citep{rei2020comet}; CometKiwi is a reference-free QE system from the WMT 2022 shared task~\citep{rei-etal-2022-cometkiwi}. Work on SI evaluation has also warned that standard MT metrics can miss interpretation-specific phenomena such as summarization, paraphrasing, disfluency, and segmentation~\citep{wein2024barriers}. In SI assessment, professional feedback can distinguish language quality, delivery/expressiveness, and promptness. We use a shared encoder with two rubric-supervised quality outputs for LQ and EXP from source text and interpreted output. The outputs remain correlated, and a one-factor sensitivity performs similarly; our experiment tests whether this quality representation complements segment-onset delay.

\paragraph{Rater-aware subjective-label modeling.}
Human labels may combine item difficulty with annotator expertise, severity, and systematic bias; interpreter-performance testing likewise documents rater severity/leniency with multifaceted Rasch analysis~\citep{han2015rater}. Crowd-label models estimate latent targets jointly with annotator reliability~\citep{whitehill2009whose,raykar2010learning}, while neural crowd layers learn annotator-specific transformations during representation learning~\citep{rodrigues2018deep}. These methods motivate evaluator-conditioned thresholds, ordinal likelihoods, and hierarchical calibration for promptness ratings. Our primary experiment uses the equal-weight aggregate target, and the rater audit characterizes evaluator-specific calibration around it.

\paragraph{Analytic professional assessment.}
Segment-level analytic assessment can preserve distinctions hidden by a single aggregate score~\citep{han2021reimagined}. This motivates our explicit rubric and our treatment of professional comments as post-hoc qualitative evidence rather than model features. Unlike work that evaluates an interpretation with an automatic metric alone, we evaluate whether automatically predicted professional quality dimensions can bridge to a separate professional promptness judgment.
